%!TeX root=../main.tex

\chapter{Metodologia}
\label{cap:metodologia}

Devido a natureza do projeto em grupo, foram adotadas diversas convenções
de documentação, controle de versão e sistemas de gerenciamento para
conservar o fluxo de desenvolvimento.

\section{Projeto}
\label{sec:projeto}

O jogo foi desenvolvido utilizando o software Godot, que é uma ferramenta
de desenvolvimento de jogos. Foi utilizada sua própria linguagem (GDScript),
que é uma linguagem de alto-nível e orientada a objetos, com uma sintaxe
similar a python. O objetivo dessa linguagem é ser otimizada e fortemente
integrada à Godot Engine, permitindo maior flexibilidade de criação e integração.
Para manter o código-fonte organizado, e legível para todos os integrantes, foi
adotada a convenção de estilo oficial do GDScript, que pode ser encontrada em sua
documentação.



\section{Desenvolvimento}
\label{sec:desenvolvimento}


\section{Arquitetura}
\label{sec:arquitetura}
